% Source PDF: source/普通高考/2012/2012山东文.pdf
% Page: 1 (q14 is a vector frequency-density histogram; pdfimages found no embedded bitmap)
% PDF embedded bitmap attempt: pdfimages -list found no image objects; page rendered at 300 dpi.
% Grid evidence: tmp/redraw_grid/2012_shandong_liberal/page1_grid100.jpg and
%   tmp/redraw_crops/2012_shandong_liberal/q14_block_grid50.jpg
% Original comparison crop: tmp/redraw_crops/2012_shandong_liberal/q14_orig.png
% Elements: x-axis break, temperature ticks 20.5--26.5 degrees C, y-axis ticks
%   0.10, 0.12, 0.18, 0.22, 0.26, six contiguous frequency-density bars, and dashed guides.
% Relations: each bar spans one 1-degree C class interval; heights are
%   (0.10, 0.12, 0.12, 0.26, 0.22, 0.18) from the source histogram.
% solid segments: x/y axes, bar outlines, x-axis break mark.
% dashed segments: horizontal guide lines at each labelled frequency-density level.
% The unused 0--20.5 range is represented by one class-width segment followed by a break.
% Each dashed guide stops at the class boundary determined by the original bar height.
\begin{tikzpicture}
\begin{axis}[
    width=9.2cm,
    height=5.2cm,
    xmin=19.5,
    xmax=28.3,
    ymin=-0.04,
    ymax=0.30,
    axis x line=middle,
    axis y line=left,
    axis line style={-stealth},
    xtick={20.5,21.5,22.5,23.5,24.5,25.5,26.5},
    xticklabels={20.5,21.5,22.5,23.5,24.5,25.5,26.5},
    x tick label as interval=false,
    ytick={0.10,0.12,0.18,0.22,0.26},
    yticklabels={,,,,},
    scaled x ticks=false,
    scaled y ticks=false,
    tick label style={font=\normalsize},
    tick align=outside,
    tick style={black},
    enlargelimits=false,
    clip=false,
]
    % Guides begin at the y-axis and stop at the corresponding bar boundary.
    \addplot[dashed,black,line width=0.35pt] coordinates {(19.5,0.10) (20.5,0.10)};
    \addplot[dashed,black,line width=0.35pt] coordinates {(19.5,0.12) (21.5,0.12)};
    \addplot[dashed,black,line width=0.35pt] coordinates {(19.5,0.18) (25.5,0.18)};
    \addplot[dashed,black,line width=0.35pt] coordinates {(19.5,0.22) (24.5,0.22)};
    \addplot[dashed,black,line width=0.35pt] coordinates {(19.5,0.26) (23.5,0.26)};

    % Histogram frequencies read from the original vector chart; class width is 1 degree C.
    \addplot[ybar interval,draw=black,fill=white,line width=0.45pt] coordinates {
        (20.5,0.10)
        (21.5,0.12)
        (22.5,0.12)
        (23.5,0.26)
        (24.5,0.22)
        (25.5,0.18)
        (26.5,0)
    };

    % The displayed 0--first-boundary segment has exactly one class width.
    \node[anchor=north east,inner sep=1pt,font=\normalsize]
        at (axis cs:19.48,-0.004) {0};
    \draw[white,line width=1.4pt]
        (axis cs:19.70,0) -- (axis cs:20.30,0);
    \draw[black,line width=0.45pt]
        (axis cs:19.70,0) -- (axis cs:19.82,0.010)
        -- (axis cs:19.94,-0.010) -- (axis cs:20.06,0.010)
        -- (axis cs:20.18,-0.010) -- (axis cs:20.30,0);

    % Preserve every ordinate label shown in the source; stagger the dense lower pair.
    \frequencyylabel[-28pt]{19.5}{0.10}{0.10}
    \frequencyylabel[-4pt]{19.5}{0.12}{0.12}
    \frequencyylabel[-4pt]{19.5}{0.18}{0.18}
    \frequencyylabel[-4pt]{19.5}{0.22}{0.22}
    \frequencyylabel[-4pt]{19.5}{0.26}{0.26}

    \node[anchor=south west,inner sep=0pt,font=\normalsize]
        at (axis cs:19.66,0.264) {\(\dfrac{\text{频率}}{\text{组距}}\)};
    \node[anchor=west,inner sep=0pt,font=\normalsize]
        at (axis cs:27.05,0.018) {平均气温/\(^{\circ}\mathrm{C}\)};
\end{axis}
\end{tikzpicture}
